\documentclass[11pt,a4paper]{beamer}
\usetheme{Madrid}
\usepackage[utf8]{inputenc}
\usepackage{amsmath}
\usepackage{amsfonts}
\usepackage{amssymb}
\usepackage{graphicx}
\usepackage{tikz}
\usepackage{booktabs}
\usepackage{array}
\usepackage{colortbl}
\usepackage{multicol}

\usefonttheme{professionalfonts}

\setbeamerfont{normal text}{size=\tiny}
\setbeamerfont{itemize/enumerate body}{size=\tiny}
\setbeamerfont{itemize/enumerate subbody}{size=\tiny}
\setbeamerfont{block body}{size=\tiny}
\setbeamerfont{block title}{size=\scriptsize}
\setbeamerfont{frametitle}{size=\small}
\setbeamerfont{title}{size=\small}
\setbeamerfont{subtitle}{size=\footnotesize}
\setbeamerfont{author}{size=\tiny}
\setbeamerfont{date}{size=\tiny}

\setlength{\leftmargini}{0.3em}
\setlength{\leftmarginii}{0.3em}
\setlength{\labelsep}{0.1em}
\setlength{\itemsep}{0.02em}
\setlength{\parskip}{0.02em}

\title[CAMS Exam Preparation]{CAMS Exam Preparation}
\subtitle{MSB, Alternative Remittance, Wash Trading \& Capital Markets}
\author{MSB Registration | Hawala | Market Manipulation | Audit}
\date{}

\setbeamercolor{title}{fg=white,bg=blue!70!black}
\setbeamercolor{frametitle}{fg=white,bg=blue!70!black}
\setbeamertemplate{navigation symbols}{}

\begin{document}

% TITLE SLIDE
\begin{frame}
\titlepage
\centering
\tiny 15 Topics | MSB | Alternative Remittance | Wash Trading | Fairness | Audit
\end{frame}

% ============================================================
% SLIDE 1 - MSB REGISTRATION OBLIGATIONS
% ============================================================
\begin{frame}{1. Money Services Business (MSB) Registration}
\tiny
\noindent\textbf{Introduction to the topic:} Money Services Businesses (MSBs) include money transmitters, check cashers, currency dealers, and prepaid access providers. Under the Bank Secrecy Act, MSBs must register with FinCEN (in the US) and with equivalent regulators in other jurisdictions. Registration is not optional — it is a legal requirement. The registration must include the MSB's name, address, ownership information, and the types of services offered. Failure to register is a criminal offense. MSBs must also maintain an AML program, designate a compliance officer, conduct customer due diligence, and file SARs.

\vspace{0.1cm}
\noindent\textbf{Type of violation or warning signs:} An entity is providing money transmission services (sending money, exchanging currency, cashing checks) but has not registered with FinCEN or the relevant regulator. The entity claims it is not an MSB because it is “only a technology platform” or “only connects buyers and sellers.” The entity has no AML program and no compliance officer. The entity processes large volumes of transactions but has never filed a SAR. The entity is owned by a person with criminal background or operates from a jurisdiction with weak AML controls.

\vspace{0.1cm}
\noindent\textbf{Correct answer approach:} The correct answer will state that MSBs must register with FinCEN (or relevant regulator) and maintain a full AML program. Eliminate answers that suggest registration is optional, that small MSBs are exempt, or that technology platforms are not MSBs. Remember: the definition of money transmission is broad — any business that accepts funds and transmits them to another location is likely an MSB.
\end{frame}

% ============================================================
% SLIDE 2 - MSB AML PROGRAM REQUIREMENTS
% ============================================================
\begin{frame}{2. MSB AML Program Requirements}
\tiny
\noindent\textbf{Introduction to the topic:} MSBs must establish and maintain a written AML program that includes: internal policies and procedures to detect and prevent money laundering; designation of a compliance officer responsible for the program; ongoing employee training; and independent testing of the program (audit). The program must be commensurate with the MSB's size and risk profile. FinCEN has issued detailed guidance on what MSB AML programs must include, including suspicious activity monitoring, currency transaction reporting, and recordkeeping.

\vspace{0.1cm}
\noindent\textbf{Type of violation or warning signs:} The MSB has no written AML program. The MSB has a compliance officer but that person has no training or authority. The MSB has never conducted independent testing of its AML program. The MSB's employees have no AML training. The MSB does not monitor transactions for suspicious activity. The MSB does not file Currency Transaction Reports (CTRs) for transactions over \$10,000. The MSB does not keep records of transactions as required by law.

\vspace{0.1cm}
\noindent\textbf{Correct answer approach:} The correct answer will state that MSBs must have a written AML program with four pillars: policies, compliance officer, training, and independent testing. Eliminate answers that suggest a smaller MSB has reduced requirements or that training is optional. Remember: the AML program must be in writing and approved by senior management. Independent testing (audit) must be conducted at least annually.
\end{frame}

% ============================================================
% SLIDE 3 - MSB CURRENCY TRANSACTION REPORTING (CTR)
% ============================================================
\begin{frame}{3. MSB Currency Transaction Reporting (CTR)}
\tiny
\noindent\textbf{Introduction to the topic:} MSBs must file Currency Transaction Reports (CTRs) with FinCEN for any transaction in currency (cash) exceeding \$10,000. This includes both single transactions and multiple transactions that appear to be structured to avoid the reporting threshold. The CTR must include the customer's name, address, identification, and the amount and nature of the transaction. MSBs that fail to file CTRs face significant penalties. Structuring — breaking transactions into amounts under \$10,000 to avoid reporting — is a criminal offense.

\vspace{0.1cm}
\noindent\textbf{Type of violation or warning signs:} A customer conducts multiple transactions just under \$10,000 on the same day at the same MSB. The customer conducts transactions just under \$10,000 at multiple MSB locations. The customer asks the MSB employee to split a transaction into smaller amounts. The MSB employee does not file a CTR because they believe the customer is “just a regular customer.” The MSB has no procedures to detect structuring. The MSB has never filed a CTR despite processing large cash volumes.

\vspace{0.1cm}
\noindent\textbf{Correct answer approach:} The correct answer will state that CTRs must be filed for cash transactions over \$10,000 and that structuring is illegal. Eliminate answers that suggest CTR filing is optional or that MSBs can aggregate transactions without reporting. Remember: the obligation to file CTRs applies to all MSBs, regardless of size. Structuring requires a SAR in addition to any CTR filing.
\end{frame}

% ============================================================
% SLIDE 4 - ALTERNATIVE REMITTANCE SYSTEMS (HAWALA)
% ============================================================
\begin{frame}{4. Alternative Remittance Systems (Hawala)}
\tiny
\noindent\textbf{Introduction to the topic:} Alternative remittance systems — also called hawala, hundi, or informal value transfer systems — are traditional money transfer methods that operate outside the formal banking system. In a hawala transaction, the sender gives cash to a hawala broker in one country, and the broker contacts a counterpart broker in another country to pay the recipient. The settlement between brokers happens later, often through goods, services, or offsetting transactions. Hawala is fast, cheap, and leaves no paper trail, making it attractive for legitimate diaspora remittances and for money laundering and terrorist financing. Many countries require hawala operators to register as MSBs, but many do not.

\vspace{0.1cm}
\noindent\textbf{Type of violation or warning signs:} A customer regularly sends money to a country through an unregistered money transfer service rather than through formal channels. The service has no office, no website, and no records. The customer cannot provide receipts or transaction confirmations. The service charges a flat fee regardless of amount, which is common in hawala. The service is located in a country with weak AML controls. The service does not collect customer identification information. The service's broker travels frequently between countries with large amounts of cash.

\vspace{0.1cm}
\noindent\textbf{Correct answer approach:} The correct answer will state that alternative remittance systems must be registered as MSBs and must comply with AML regulations, including CDD and SAR filing. Eliminate answers that suggest hawala is exempt from AML rules because it is traditional or cultural. Remember: the FATF explicitly includes informal value transfer systems in its recommendations — they are not exempt.
\end{frame}

% ============================================================
% SLIDE 5 - HAWALA RED FLAGS AND DETECTION
% ============================================================
\begin{frame}{5. Hawala Red Flags and Detection}
\tiny
\noindent\textbf{Introduction to the topic:} Detecting hawala transactions is challenging because they leave no paper trail. However, red flags can be observed: customers who send money to high-risk jurisdictions through informal channels; customers who receive frequent international wires from countries with no legitimate business relationship; customers who ask about money transfer services that do not require identification; customers who send money in round amounts (e.g., \$5,000, \$10,000) with no business purpose; and customers who send money to individuals in countries where they have no family or business ties.

\vspace{0.1cm}
\noindent\textbf{Type of violation or warning signs:} Customer sends regular remittances to a country through an unregistered service. Customer refuses to use formal banks or MSBs, claiming they are “too expensive” or “too slow.” Customer's declared income cannot explain the volume of remittances. Customer sends money to countries known for high terrorism risk or money laundering. Customer receives funds from a hawala broker and then deposits them into a bank account. Customer has no receipts or records of the transactions.

\vspace{0.1cm}
\noindent\textbf{Correct answer approach:} The correct answer will state that banks and MSBs must identify customers who use alternative remittance systems and file SARs when hawala activity is suspected. Eliminate answers that suggest hawala transactions cannot be detected or that banks have no responsibility to inquire about the method of transfer. Remember: when a customer uses informal channels, it suggests they are trying to avoid scrutiny.
\end{frame}

% ============================================================
% SLIDE 6 - WASH TRADING DEFINITION AND RECOGNITION
% ============================================================
\begin{frame}{6. Wash Trading Definition and Recognition}
\tiny
\noindent\textbf{Introduction to the topic:} Wash trading is the practice of buying and selling the same security at the same price, creating artificial trading volume with no change in beneficial ownership. The transactions cancel each other out, but they create the appearance of active trading. Wash trading is illegal in most jurisdictions because it manipulates market prices and creates false volume that can attract other investors. It is also a predicate offense for money laundering — the artificial activity can be used to make illicit funds appear to come from legitimate trading.

\vspace{0.1cm}
\noindent\textbf{Type of violation or warning signs:} The same customer buys and sells the same security within a short period at the same price. Two accounts with common control (e.g., same beneficial owner) buy and sell the same security to each other. The customer executes trades that have no economic purpose — no price movement, no profit, no loss. The customer's trading volume is unusually high relative to their net worth. The customer asks the broker to “match” trades at specific prices. The customer trades the same security repeatedly with no change in position.

\vspace{0.1cm}
\noindent\textbf{Correct answer approach:} The correct answer will state that wash trading is market manipulation and a money laundering predicate offense requiring SAR filing. Eliminate answers that suggest wash trading is only an exchange rule violation or that it is acceptable in certain circumstances. Remember: the key indicator is the absence of economic purpose — if the trades do not change the customer's position, they are likely wash trades.
\end{frame}

% ============================================================
% SLIDE 7 - WASH TRADING IN CRYPTOCURRENCY MARKETS
% ============================================================
\begin{frame}{7. Wash Trading in Cryptocurrency Markets}
\tiny
\noindent\textbf{Introduction to the topic:} Wash trading is particularly common in cryptocurrency markets, where exchanges may inflate trading volume to attract users. A single person or entity can control multiple accounts on the same exchange and trade between them to create volume. This practice creates fake liquidity and can manipulate prices. Crypto wash trading is illegal in regulated jurisdictions and is a predicate offense for money laundering. The illicit proceeds from wash trading may be used to launder funds by creating a false appearance of legitimate trading activity.

\vspace{0.1cm}
\noindent\textbf{Type of violation or warning signs:} A cryptocurrency exchange reports unusually high trading volume but has few users. A customer trades the same cryptocurrency back and forth between accounts they control. The customer's trades are at the same price with no price movement. The customer executes trades that generate no profit or consistent small losses. The customer uses multiple accounts on the same exchange with the same IP address or device. The customer's wash trading coincides with deposits of large amounts of cryptocurrency from unknown sources.

\vspace{0.1cm}
\noindent\textbf{Correct answer approach:} The correct answer will state that wash trading in crypto is market manipulation and money laundering, requiring SAR filing and exchange enforcement. Eliminate answers that suggest crypto wash trading is not a violation because it is “decentralized” or “not regulated.” Remember: regulated exchanges have AML obligations, and wash trading on any exchange is illegal if it manipulates the market.
\end{frame}

% ============================================================
% SLIDE 8 - CAPITAL MARKET FAIRNESS (INSIDER TRADING)
% ============================================================
\begin{frame}{8. Capital Market Fairness - Insider Trading}
\tiny
\noindent\textbf{Introduction to the topic:} Capital market fairness requires that all investors have equal access to material information. Insider trading violates this principle by allowing individuals with non-public information to profit unfairly. Insider trading is illegal in most jurisdictions and is a predicate offense for money laundering. The profits from insider trading are considered criminal proceeds. Broker-dealers and investment banks must have surveillance systems to detect suspicious trading patterns that indicate insider trading.

\vspace{0.1cm}
\noindent\textbf{Type of violation or warning signs:} A customer with no history of profitable trading suddenly makes highly profitable trades just before a major corporate announcement. The customer trades in a stock that is not typically traded by that customer. The customer cannot explain the source of their trading information. The customer trades in the same stock as another account controlled by a person with access to inside information. The customer closes all positions immediately after the announcement, regardless of price. The customer uses a complex structure (trust, shell company, offshore account) to hide ownership.

\vspace{0.1cm}
\noindent\textbf{Correct answer approach:} The correct answer will state that insider trading is a predicate offense for money laundering and requires SAR filing, and that broker-dealers must monitor for suspicious trading patterns. Eliminate answers that suggest insider trading is only a securities violation and not an AML concern. Remember: the profits from insider trading are criminal proceeds — they must be reported.
\end{frame}

% ============================================================
% SLIDE 9 - CAPITAL MARKET FAIRNESS (MANIPULATION)
% ============================================================
\begin{frame}{9. Capital Market Fairness - Market Manipulation}
\tiny
\noindent\textbf{Introduction to the topic:} Market manipulation includes practices that artificially affect the price of a security, such as pump-and-dump schemes, spoofing (placing and canceling orders to create false demand), and layering (placing multiple orders at different prices that are canceled before execution). These practices undermine capital market fairness and are predicate offenses for money laundering. The proceeds from manipulation are criminal proceeds. Broker-dealers must have surveillance systems to detect these patterns.

\vspace{0.1cm}
\noindent\textbf{Type of violation or warning signs:} A customer places large orders and cancels them immediately before execution (spoofing). A customer places multiple orders at different prices and cancels them after the market moves (layering). A customer trades a stock with no legitimate business reason, causing the price to move. A customer sells large blocks of a stock after the price increases without a news catalyst. A customer engages in pump-and-dump activity — hyping the stock and then selling. The customer uses multiple accounts to manipulate the price.

\vspace{0.1cm}
\noindent\textbf{Correct answer approach:} The correct answer will state that market manipulation is a predicate offense for money laundering and requires SAR filing, and that broker-dealers must have surveillance systems to detect spoofing, layering, and pump-and-dump patterns. Eliminate answers that suggest market manipulation is only a regulatory issue and not an AML concern. Remember: manipulation profits are criminal proceeds.
\end{frame}

% ============================================================
% SLIDE 10 - AUDIT AND INDEPENDENT TESTING
% ============================================================
\begin{frame}{10. Audit and Independent Testing of AML Programs}
\tiny
\noindent\textbf{Introduction to the topic:} AML programs must be independently tested (audited) on a regular basis. For MSBs, the audit must be conducted at least annually. For banks and broker-dealers, the audit frequency depends on the institution's risk profile, but at least annually is typical. The independent testing must be conducted by a qualified person (internal or external) who is independent of the AML function. The audit should assess the effectiveness of the AML program, identify weaknesses, and make recommendations for improvement. The findings must be reported to senior management and the board of directors.

\vspace{0.1cm}
\noindent\textbf{Type of violation or warning signs:} The institution has not conducted an AML audit in over a year. The institution's audit is conducted by the same person who manages the AML program (no independence). The audit was conducted but the findings were not presented to the board. The audit identified significant weaknesses but no corrective action was taken. The institution has no audit schedule or plan. The audit was a “check-the-box” exercise with no meaningful analysis. The auditors lacked relevant AML expertise.

\vspace{0.1cm}
\noindent\textbf{Correct answer approach:} The correct answer will state that AML programs require independent testing at least annually, by qualified personnel, with findings reported to senior management and the board. Eliminate answers that suggest audit is optional or that internal staff (without independence) can conduct the audit. Remember: independence is the key requirement — the auditor cannot report to the compliance officer.
\end{frame}

% ============================================================
% SLIDE 11 - AUDIT FINDINGS AND CORRECTIVE ACTION
% ============================================================
\begin{frame}{11. Audit Findings and Corrective Action}
\tiny
\noindent\textbf{Introduction to the topic:} AML audits often identify deficiencies such as gaps in customer due diligence, inadequate transaction monitoring, incomplete SAR documentation, or insufficient employee training. When deficiencies are identified, the institution must implement corrective actions within a reasonable timeframe. The board of directors must oversee the corrective action plan. Failure to implement corrective actions can result in regulatory enforcement actions, fines, and in severe cases, criminal prosecution.

\vspace{0.1cm}
\noindent\textbf{Type of violation or warning signs:} Audit findings identified significant AML deficiencies but management did not take corrective action. The institution has a backlog of unresolved audit findings from previous years. The board was informed of deficiencies but did not take action. The corrective actions implemented were inadequate or did not address the root cause. The institution's AML program remained deficient for an extended period despite audit findings. The auditor found that the same deficiencies existed in consecutive audits.

\vspace{0.1cm}
\noindent\textbf{Correct answer approach:} The correct answer will state that audit findings require timely corrective action, with board oversight, and failure to implement corrective actions is a regulatory violation. Eliminate answers that suggest audit findings are merely recommendations or that corrective action is optional. Remember: unresolved audit findings are a serious red flag for regulators.
\end{frame}

% ============================================================
% SLIDE 12 - BOARD OVERSIGHT OF AML PROGRAM
% ============================================================
\begin{frame}{12. Board Oversight of AML Program}
\tiny
\noindent\textbf{Introduction to the topic:} The board of directors is ultimately responsible for the AML program. The board must approve the AML program, receive regular reports on its effectiveness, and ensure adequate resources are allocated to AML compliance. The board must also be informed of significant AML findings, SAR filings, and regulatory actions. Failure of board oversight is a culture failure and has been the basis for enforcement actions against large banks and broker-dealers.

\vspace{0.1cm}
\noindent\textbf{Type of violation or warning signs:} The board has never reviewed the AML program. The board receives AML reports but does not discuss them or ask questions. The board has not approved the AML program in writing. The board does not receive SAR statistics or summaries. The board is not informed of AML audit findings. The board is not informed of regulatory examinations or findings. The board does not ensure adequate AML resources (staff, technology, budget).

\vspace{0.1cm}
\noindent\textbf{Correct answer approach:} The correct answer will state that the board of directors has ultimate responsibility for the AML program and must approve it, receive reports, and ensure resources. Eliminate answers that suggest AML oversight is solely a compliance function or that the board's role is limited. Remember: tone from the top starts with the board. If the board is not engaged, the AML program is likely inadequate.
\end{frame}

% ============================================================
% SLIDE 13 - COMPLIANCE CULTURE IN CAPITAL MARKETS
% ============================================================
\begin{frame}{13. Compliance Culture in Capital Markets}
\tiny
\noindent\textbf{Introduction to the topic:} A strong culture of compliance is essential in capital markets, where pressures to generate revenue can conflict with AML obligations. Traders and relationship managers may be incentivized to ignore red flags if their compensation is tied to trading volume or revenue. Firms must ensure that compliance is not sacrificed for revenue. This includes risk-adjusted compensation, independent compliance authority, and a board that actively oversees AML.

\vspace{0.1cm}
\noindent\textbf{Type of violation or warning signs:} Traders receive bonuses based solely on trading volume with no risk adjustment. Compliance recommendations are overridden by senior management citing revenue impact. The firm's largest customers are its highest-risk customers, and the firm has never filed SARs on them. The board has never reviewed AML metrics. Employees report that compliance is seen as a “barrier” to business. The firm has had multiple regulatory enforcement actions for AML violations.

\vspace{0.1cm}
\noindent\textbf{Correct answer approach:} The correct answer will state that a strong culture of compliance requires risk-adjusted compensation, independent compliance authority, and board oversight, and that revenue cannot override AML obligations. Eliminate answers that suggest compliance is secondary to revenue or that traders have no AML obligations. Remember: culture failures are increasingly the basis for enforcement actions and personal liability for senior management.
\end{frame}

% ============================================================
% SLIDE 14 - SURVEILLANCE SYSTEMS FOR CAPITAL MARKETS
% ============================================================
\begin{frame}{14. Surveillance Systems for Capital Markets}
\tiny
\noindent\textbf{Introduction to the topic:} Broker-dealers and investment banks must have trade surveillance systems to detect insider trading, market manipulation, and other suspicious activities. These systems monitor trading patterns, order book activity, and communications. Surveillance systems must be calibrated to the firm's risk profile and cover all products traded. Firms must have procedures for investigating alerts and filing SARs when suspicious activity is confirmed.

\vspace{0.1cm}
\noindent\textbf{Type of violation or warning signs:} The firm has no trade surveillance system. The surveillance system generates alerts but no one investigates them. The surveillance system is not calibrated to detect the firm's specific products or trading patterns. The firm has never filed a SAR from a surveillance alert. The firm's surveillance system does not cover all products (e.g., options, futures, fixed income). The firm relies only on manual reviews with no automated system. The firm has no procedure for escalating surveillance findings to compliance.

\vspace{0.1cm}
\noindent\textbf{Correct answer approach:} The correct answer will state that broker-dealers must have trade surveillance systems that are calibrated to their risk profile, with procedures for investigating alerts and filing SARs. Eliminate answers that suggest manual review is sufficient or that surveillance is optional. Remember: surveillance must cover all products and all trading desks. Uninvestigated alerts are as bad as no alerts.
\end{frame}

% ============================================================
% SLIDE 15 - MSB EXAMINATION AND ENFORCEMENT
% ============================================================
\begin{frame}{15. MSB Examination and Enforcement}
\tiny
\noindent\textbf{Introduction to the topic:} MSBs are subject to examination by FinCEN and state regulators. Examinations assess the MSB's AML program, customer due diligence, transaction monitoring, CTR and SAR filing, and independent testing (audit). Enforcement actions can include civil penalties (fines), cease and desist orders, and criminal prosecution. MSBs that fail to register or fail to maintain an adequate AML program face significant penalties. Individuals who run unregistered MSBs can face criminal charges.

\vspace{0.1cm}
\noindent\textbf{Type of violation or warning signs:} The MSB has never been examined by a regulator. The MSB has no record of regulatory correspondence. The MSB has no AML program or an inadequate program. The MSB has no compliance officer. The MSB has never filed a CTR or SAR despite processing large volumes. The MSB's owner has a criminal background. The MSB operates from a residential address or a location with no business signage. The MSB processes transactions for customers from high-risk jurisdictions.

\vspace{0.1cm}
\noindent\textbf{Correct answer approach:} The correct answer will state that MSBs are subject to examination and enforcement and must maintain a full AML program, file CTRs and SARs, and register with FinCEN. Eliminate answers that suggest smaller MSBs are exempt from examination or that the owner's background is not relevant. Remember: unregistered MSB operations are a criminal offense. Exam questions often include scenarios where an entity is operating as an MSB without registration — the correct answer is that registration is required and failure to register is illegal.
\end{frame}

% ============================================================
% SUMMARY SLIDE - MSB, ALTERNATIVE REMITTANCE, WASH TRADING
% ============================================================
\begin{frame}{Summary: Key Principles for MSB, Remittance, and Capital Markets}
\tiny
\noindent\textbf{How to answer exam questions on these topics:}

\vspace{0.1cm}
\noindent• \textbf{MSB Registration:} Any entity transmitting money must register as an MSB. No exceptions for technology platforms or informal systems.

\vspace{0.1cm}
\noindent• \textbf{MSB AML Program:} Must have written policies, compliance officer, training, and independent testing. All required.

\vspace{0.1cm}
\noindent• \textbf{CTR Filing:} Cash transactions over \$10,000 require CTRs. Structuring to avoid CTRs is illegal.

\vspace{0.1cm}
\noindent• \textbf{Alternative Remittance (Hawala):} Must be registered as an MSB. Not exempt because it is traditional or cultural.

\vspace{0.1cm}
\noindent• \textbf{Wash Trading:} Market manipulation and money laundering predicate offense. No economic purpose = red flag.

\vspace{0.1cm}
\noindent• \textbf{Insider Trading:} Predicate offense for money laundering. Unusual profits before announcements = red flag.

\vspace{0.1cm}
\noindent• \textbf{Market Manipulation:} Spoofing, layering, pump-and-dump = market manipulation and money laundering.

\vspace{0.1cm}
\noindent• \textbf{AML Audit:} Independent testing required at least annually. Findings must be reported to board.

\vspace{0.1cm}
\noindent• \textbf{Board Oversight:} Board must approve AML program, receive reports, and ensure resources. Ultimate responsibility.

\vspace{0.1cm}
\noindent• \textbf{Compliance Culture:} Revenue cannot override AML. Risk-adjusted compensation required.

\vspace{0.1cm}
\noindent• \textbf{Surveillance:} Must cover all products, investigate alerts, and file SARs. Uninvestigated alerts = deficiency.

\vspace{0.1cm}
\noindent• \textbf{Examination:} MSBs are subject to examination and enforcement. Failure to register or maintain program = penalties.

\vspace{0.3cm}
\centering
{When in doubt: Register, report, audit, and file SAR. MSBs and capital markets are heavily regulated — compliance is not optional.}
\end{frame}

\end{document}